 \chapter{Methodology}

\section{Experimental Design}
    
To investigate the influence of the construction strategy on the quality of the BVH, an experiment is implemented. 
Algorithms from three different construction classes are to be implemented: top-down, bottom-up, and incremental, following the classification presented by Ericson\cite[p.~239]{ericson2004real} % Seitenzahl fehlt hier noch
The experiment will focus on static two-dimensional collision detection, therefore particles will not move after dataset generation, and each hierarchy is generated completely before traversal is done.

The primary independent variable of the experiment is the BVH construction strategy. 
Particle count and scene coverage are varied as additional experimental factors. 
The primary dependent variables are the number of AABB overlap tests performed during traversal.
Each comparison bewteen AABBs is counted as a bounding-volume check, including those between two leaf AABBs.
Additionally, the number of primitive checks is recorded as a control metric, since identical leaf AABBs should always produce identical candidate pairs. % Das ist streng genommen ein Finding meiner Arbeit, ich wusste es ja vorher nicht, trotzdem führe ich es nicht mehr als primäre Metrik und schließe es auch aus der Forschungsfrage aus. Ich denke das ist sinnvoll da sie für die Discussion sowieso basically übergangen werden

To isolate the effect of the construction strategy, all other components of the collision detection procedure are kept constant between experiment runs.
All strategies receive the same datasets and use the same geometric representation, collision detection and traversal procedure. 
This is meant to ensure, that differences in the recorded number of checks can primarily be attributed to the differences in the structure generated by the respective construction algorithms.

The experiment uses particle counts of 25, 50, 100, 200, 500 and 1000, coverage values of 0.05, 0.10 and 0.20, a particle radius of 1, and 50 datasets per parameter combination generated from a base seed of 42.
All three algorithms are evaluated on the same datasets.

\section{Dataset generation}
Each dataset contains a number of equally sized circular particles, contained in a square two-dimensional area.
The particle centers are taken from a uniform distribution using Python's \texttt{random} function and a deterministic seed generated from the experiment's parameters.
Both x and y coordinates are limited to $[r,L-r]$, where r is the radius of the particles, ensuring they are completely enclosed in the bounds of the square.

To keep the nominal particle coverage approximately constant for different particle counts, a target coverage (c) is defined, and the domain side lengths are calculated as:

\[
L=\sqrt{\frac{n\pi r^2}{c}}.
\]

Keeping the particle coverage constant keeps the particle density approximately comparable across different particle counts.
As particles may overlap, the actual occupied area may differ from the nominal target.


\section{BVH construction}
The general principles of the three construction mechanisms are explained in the literature review section.
This section therefore focuses on the specific implementations used in this experiment.
The results primarily apply to these implementations rather than to the respective algorithm classes as a whole.

\subsection{Top-down Construction}

The top-down implementation uses a median split heuristic. 
For each node its AABB is compared along both axes. 
The particles are then sorted by their center coordinate along the longer axis and split at the median position.

\subsection{Bottom-up Construction}

The bottom-up implementation uses the area of the merged AABB as its criterion for clustering. 
In each construction step, all possible pairs of still active nodes are evaluated and the pair that produces the smallest merged AABB is selected.
The selected pair is then replaced by the newly created parent node. This process is repeated until only the root remains.

\subsection{Incremental Construction}

Particles are inserted in their generated dataset order.
For each new particle the tree is traversed, by following the child whose AABB requires the smallest area increase to include the new particle.
The nodes along this search path are then compared using an insertion-cost function, that accounts for the area of the AABB enclosing the candidate node and the new particle, as well as resulting area increases of its ancestor. 
The node with the lowest cost is selected as the position for insertion.

\section{Traversal and Evaluation}

After construction, each hierarchy is queried using the same self-traversal procedure. 
A pair of subtrees is discarded when their AABBs do not overlap. 
If they do overlap, traversal continues with the remaining child-node combinations. 
A primitive test, which is defined as \(\lVert p_i-p_j\rVert \le r_i+r_j\), where p are the center coordinates of the particles, and r is the radius of the circles, is performed once two leaf nodes with overlapping AABBs are reached.
The self-traversal compares sibling subtrees recursively so that each relevant unordered leaf pair is reached once.

Every evaluation records the number of necessary bounding volume overlap tests, and the number of primitive collision tests separately.
Keeping the measures separately was chosen, as they represent two different operations that potentially come with a different computational cost, which therefore is relevant for measuring tree-quality.

Tree height is the maximum edge depth with the root at zero. 
Total internal AABB area sums all non-leaf AABBs. 
Both are supplementary rather than direct measurements of traversal work.

Each generated datasets structure is validated, and its candidate and collision pairs are compared against brute-force results as a reference.




\newpage